\documentclass[12pt]{article}

\usepackage[utf8]{inputenc}
\usepackage{amsmath, amssymb, amsthm}
\usepackage{geometry}
\geometry{a4paper, margin=2.5cm}

\title{Reprezentarea Standard și Unicitatea Factorilor Invarianti}
\author{}
\date{}

\theoremstyle{definition}
\newtheorem{definition}{Definiție}

\theoremstyle{plain}
\newtheorem{theorem}{Teoremă}

\begin{document}

\maketitle

\begin{definition}
Fie $R$ un domeniu principal (PID). Un $R$-modul finit generat $M$ se numește
\emph{reprezentat în formă standard} dacă există elemente nenule


\[
d_1, \dots, d_m \in R, \qquad d_1 \mid d_2 \mid \cdots \mid d_m,
\]


și un întreg $k \ge 0$ astfel încât


\[
M \cong R^k \;\oplus\; R/d_1 R \;\oplus\; R/d_2 R \;\oplus\; \cdots \;\oplus\; R/d_m R.
\]


Numerele $d_1, \dots, d_m$ se numesc \emph{factori invarianti} ai lui $M$.
\end{definition}

\begin{theorem}[Reprezentarea standard]
Fie $R$ un PID și $M$ un $R$-modul finit generat. Atunci $M$ admite o reprezentare standard, adică există


\[
k \ge 0,\quad d_1, \dots, d_m \in R \setminus \{0\},\quad d_1 \mid d_2 \mid \cdots \mid d_m,
\]


astfel încât


\[
M \cong R^k \oplus R/d_1 R \oplus \cdots \oplus R/d_m R.
\]


Mai mult, putem alege elemente $x_1, \dots, x_m \in M$ astfel încât


\[
\operatorname{ann}(x_i) = d_i R.
\]


\end{theorem}

\begin{theorem}[Unicitatea reprezentării standard]
În ipotezele teoremei precedente, numerele $k$ și $d_1, \dots, d_m$ sunt unice
până la unități din $R$.

Mai precis:
\begin{enumerate}
    \item[(i)] rangul liber $k$ este unic;
    \item[(ii)] numărul $m$ de factori torsionali este unic;
    \item[(iii)] fiecare $d_i$ este unic până la înmulțire cu o unitate din $R$;
    \item[(iv)] condiția de divizibilitate
    

\[
    d_1 \mid d_2 \mid \cdots \mid d_m
    \]


    determină în mod unic șirul factorilor invarianti.
\end{enumerate}
\end{theorem}

\begin{theorem}[Caracterizare prin anihilatori]
În reprezentarea standard, elementele $x_1, \dots, x_m$ pot fi alese astfel încât


\[
M_{\text{tors}} = Rx_1 \oplus \cdots \oplus Rx_m,
\]


iar


\[
\operatorname{ann}(x_i) = d_i R.
\]


Astfel, factorii invarianti sunt determinați de anihilatorii generatorilor torsionali.
\end{theorem}

\end{document}
